\scanpage{95}
\margin{Dualité pour catégories fibrées et cofibrées}

Catégorie

\diagram{\begin{tikzcd}[ampersand replacement=\&] \unsure{E|\mathcal E}\arrow[d,"p"'] \\ B \end{tikzcd}}

sur une autre

1) La catégorie est \uline{fibrée} ou $p$ fibrant, si pour \deleted{\unsure{toute}} flèche $f:x\longrightarrow y$ de $B$, on peut définir $f^*:E_y\mathrel{\rule{1.4em}{.35pt}}E_x$, avec transitivité \unsure{essentielle|essent\textsuperscript{le}} :
\inserted{$\operatorname{Hom}_f(X,Y)\simeq\operatorname{Hom}_{\mathrm{id}_x}(X,f^*(Y))$}
On trouve un pseudo-foncteur

\[ F:B^\circ\longrightarrow(\mathrm{Cat}) \]

2) La catégorie est \uline{cofibrée}, ou $p$ cofibrant, si $p^\circ$ est fibrant, i.e. si pour toute $f:x\mathrel{\rule{1.4em}{.35pt}}y$ de $B$, existe un foncteur $f_*:E_x\mathrel{\rule{1.4em}{.35pt}}E_y$, avec transitivités essentielles : $\operatorname{Hom}_f(X,Y)\simeq\operatorname{Hom}(f_*(\unsure{x|X|X'}),Y)\rbrace$. Il donne naissance à un pseudo-foncteur

\[ F:B\longrightarrow(\mathrm{Cat}) \]

\uline{\unsure{Liens|Lien}} a) Si $p$ est fibrant, pour que (pour $f:x\mathrel{\rule{1.4em}{.35pt}}y$ donnée) $f_*$ existe, il f. et s. que $f^*$ ait un adjoint à g. Pour que $p$ soit cofibrant, il f. et s. que les $f^*$ aient tous des adjoints à g.

b) Si $p$ est cofibrant, pour que $f^*$ existe, il f. et s. que $f_*$ ait un adjoint à droite. Pour que $p$ soit fibrant, il f. et s. que tous les $f_*$ aient des adjoints à droite.

La catégorie \uline{opposée relative} d'une catégorie fibrée $p:E\mathrel{\rule{1.4em}{.35pt}}B$ se définit seulement à équivalence près, c'est une autre catégorie fibrée, associée au pseudo-foncteur

% archive sheet number: 94
